\documentclass[10pt,a4paper]{amsart}
\usepackage[margin=19mm]{geometry}
\usepackage{amsmath,amssymb,mathtools}
\usepackage[scr=rsfs,cal=euler]{mathalfa}
\usepackage{xfrac}
\usepackage[unicode,hidelinks,bookmarks=false]{hyperref}
\newcommand{\ie}{i.e.\ }
\newcommand{\cf}{cf.\ }
\newcommand{\resp}{respectively\ }
\newcommand{\Subsection}{\subsection}
\providecommand{\nobreakdash}{\nobreak\hbox{-}}
\setlength{\emergencystretch}{2em}
\multlinegap0pt
\usepackage{tikz}
\usepackage{amssymb}
\usepackage{xcolor}
\newtheorem{thm}{Theorem}
\title{{The second and third Hankel determinants for certain classes of functions}}
\author{Vasudevarao Allu, Shobhit Kumar}
\date{Source: arXiv:2604.10298v1 (CC BY 4.0)}
\begin{document}
\maketitle

\noindent\emph{Source and licence.} {The second and third Hankel determinants for certain classes of functions}. Version arXiv:2604.10298v1 (CC BY 4.0), distributed by arXiv under the Creative Commons Attribution 4.0 International licence, http://creativecommons.org/licenses/by/4.0/, as recorded in the arXiv licence metadata for this exact version and shown on the abstract page https://arxiv.org/abs/2604.10298v1. Full licence notice: \texttt{licenses/arxiv-2604.10298-CC-BY-4.0.txt}.\par\smallskip

\noindent\textit{Editorial scope.} Counted unit: the article's thm thm:H3-1 (source0.tex lines 1064--1070). The article's complete original proof of it is delivered with it (source0.tex lines 1071--1866, one \texttt{proof} environment of 796 lines). No mathematics is written, repaired or re-derived: the statement and the proof are the article's own lines, in the article's own notation, and no retained line is substituted or re-flowed.  No further source-local context is needed: every cross-reference of every retained line resolves inside the two delivered spans.  Nothing was omitted from any retained span, so the package carries no omission record.  No retained line contains a citation, so the wrapper carries no bibliography and no bibliography entry.  No retained line contains a figure environment, an image-inclusion command or drawing code, so no image is delivered and the package declares no asset dependency.  Editorial formatting only, in addition: an AMS \texttt{amsart} preamble that restates the article's own declarations and macros used by the retained lines, namely the environments \texttt{thm} on separate counters, together with the standard packages those lines need.  The retained environments print in this document as Theorem 1 in order of appearance, whereas the article prints the same items with the numbering of its own full document.  The complete native source of the article as distributed in the arXiv e-print for this exact version is bundled unchanged as \texttt{sources/198259/source0.tex}.
\begin{thm}\label{thm:H3-1}
Let $f\in \mathcal{S}^*(\varphi)$. Then
$
|H_3(1)|\leq {1}/{9},
$
where $H_3(1)$ denote the third Hankel determinant. Moreover, the estimate is sharp.
\end{thm}

\begin{proof}
Let $f\in \mathcal{S}^*(\varphi),$ be given by 
\begin{align*}
	\frac{z f'(z)}{f(z)}\prec \varphi(z), \quad z\in \mathbb{D}.
\end{align*}
 The third Hankel determinant is defined by:
\[
H_3(1)=
\begin{vmatrix}
	1 & a_2 & a_3 \\
	a_2 & a_3 & a_4 \\
	a_3 & a_4 & a_5
\end{vmatrix}
= a_3(a_2 a_4 - a_3^2)-a_4(a_4-a_2a_3)+a_5(a_3-a_2^2).
\]
For $f\in\mathcal{S}^*(\varphi)$ there exists a Schwarz function $w$ such that
\begin{align}\label{eq:4.4}
	\frac{z f'(z)}{f(z)}=\varphi(w(z))=\left(1+\frac{w(z)}{2}\right)^2.
\end{align}
We derive the coefficients with respect to the Schwarz coefficients.
Writing the Schwarz function in the series form such that
\[
w(z)=\sum_{n=1}^{\infty} c_n z^n,
\]
and putting the expansions of $f\in \mathcal{S}^*(\varphi)$ and $w(z)$ in \eqref{eq:4.4},
we obtain the coefficient relations
\begin{equation}
	\left.
	\begin{aligned}
		a_2 &= c_1,\\
		a_3 &= \frac{1}{8}\left(5c_1^2+4c_2\right),\\
		a_4 &= \frac{1}{24}\left(7c_1^3+16c_1c_2+8c_3\right),\\
		a_5 &= \frac{1}{384}\left(43c_1^4+184c_1^2c_2+72c_2^2+176c_1c_3+96c_4\right).
	\end{aligned}
	\right\}
\end{equation}
Substituting these values into the expansion of $H_3(1)$ yields
\begin{align}\label{eq:5.2}
	9216\, H_{3}(1)=&
	-61c_1^6 + 244c_1^4c_2 + 464c_1^3c_3 + 1088c_1c_2c_3
	- 8c_1^2(89c_2^2 + 108c_4)\notag \\
	\qquad
	&- 32(9c_2^3 + 32c_3^2 - 36c_2c_4).
\end{align}
We now use the relations between the coefficients of the Schwarz functions, which can be stated as:
\begin{align}\label{eq:5.3}
	c_2 &= (1-c_1^2)\gamma,\notag\\
	c_3 &= (1-c_1^2)\left((1-|\gamma|^2)\eta - c_1\gamma^2\right),\notag\\
	c_4 &= (1-c_1^2)\left(c_1^2\gamma^3 - (1-|\gamma|^2)\left(2c_1\gamma\eta + \overline{\gamma}\eta^2 - (1-|\eta|^2)\rho\right)\right),
\end{align}
where $|\gamma|,\,  |\eta|,\,  |\rho| \in [0, 1].$\\[2mm]
Using coefficient relation in \eqref{eq:5.3} in \eqref{eq:5.2}, we obtain
\begin{equation}\label{eq:5.4}
	9216 \, H_{3}(1)=A_{1}+B_{1} \eta +C_{1}\eta^2+D_{1}\rho
\end{equation}
where,
\begin{align*}
	A_1=&-61c_1^6-244\gamma c_1^4(-1+c_1^2)+128\gamma^4c_1^2(-1+c_1^2)^2 \\
	\quad &-8\gamma^2c_1^2\left(89-120c_1^2+31c_1^4\right)
	+32\gamma^3\left(-9-7c_1^2+14c_1^4+2c_1^6\right),\\[0.5em]
	B_1=&16\left(-1+|\gamma|^2\right)c_1(-1+c_1^2)
	\left(29c_1^2+16\gamma^2(-1+c_1^2)+\gamma(68+40c_1^2)\right),\\[0.5em]
	C_1=&-32\left(-1+|\gamma|^2\right)(-1+c_1^2)
	\Bigl(-32(-1+c_1^2)+32|\gamma|^2(-1+c_1^2) \\
	\quad &-9\left(3c_1^2+4\gamma(-1+c_1^2)\right)\overline{\gamma}\Bigr),\\[0.5em]
	D_1=&288\left(-1+|\gamma|^2\right)\left(-1+|\eta|^2\right)(-1+c_1^2)
	\left(3c_1^2+4\gamma(-1+c_1^2)\right).
\end{align*}
Since the class $\mathcal{S}^*(\varphi)$ is invariant under rotations, we may assume without loss of generality that $c_1$ is real and nonnegative. Thus
$
0\le c_1\le 1.
$
Setting
\[
c_1=p_1,\qquad x=|\gamma|,\qquad y=|\eta|,
\]
and taking modulus on both sides of \eqref{eq:5.4}, using the triangle inequality and $|\rho|\le 1$, we obtain
\[
9216\,|H_3(1)|\le H(p_1,x,y),
\]
where 
\begin{align}
	H(p_1,x,y)=&61p_1^6+244p_1^4(1-p_1^2)x
	+8p_1^2(89-120p_1^2+31p_1^4)x^2 \notag\\
	\quad
	&-32(-9-7p_1^2+14p_1^4+2p_1^6)x^3
	+128p_1^2(-1+p_1^2)^2x^4 \notag\\
	\quad
	&+16(1-x^2)p_1(1-p_1^2)\bigl(29p_1^2+16x^2(1-p_1^2)+x(68+40p_1^2)\bigr)y \notag\\
	\quad
	&+32(1-x^2)(1-p_1^2)\bigl(32(1-x^2)(1-p_1^2)+9(3p_1^2+4x(1-p_1^2))x\bigr)y^2 \notag\\
	\quad
	&+288(1-x^2)(1-y^2)(1-p_1^2)\bigl(3p_1^2+4x(1-p_1^2)\bigr).
\end{align}
Now our aim is to maximize this quantity when $\mathcal{D}:=\{(p_1, x, y)\in \mathbb{R}^3:(p_1,x,y)\in [0,1]\times[0, 1]\times[0, 1]\}.$
It is indeed easy to check that the coefficient of $y$ remains non-negative over the region $\mathcal{D}.$ Therefore we take $y=1$ for the specific term 
\[
16(1-x^2)p_1(1-p_1^2)\bigl(29p_1^2+16x^2(1-p_1^2)+x(68+40p_1^2)\bigr)y
\]
and denote this new function as $H_{1}(p_1, x, y),$ with $H(p_1, x, y)\le H_{1}(p_1, x, y).$ Since it is a quadratic polynomial in the variable $y,$ 
\[
\max_{\mathcal{D}}H_{1}(p_1, x, y)=\max_{(p_{1},x)\in [0, 1]\times[0, 1]}\left\{H_{1}(p_1, x, 0), H_{1}(p_1, x, 1)\right\}.
\]
Let us denote
\[
G_{1}(p_1, x):= H_{1}(p_1, x, 1)
\qquad\text{and}\qquad
G_{2}(p_1, x):= H_{1}(p_1, x, 0).
\]
We first deal with the function $G_{1}(p_1, x).$ In this case we prove that the function attains its maximum at the point $(0, 0),$ which is $G_{1}(0, 0)=1024$.\\[2mm]
Let
\[
F(p_1,x):=1024-G_1(p_1,x),
\qquad p:=p_1.
\]
We divide the original square \([0,1]^2\) into the four rectangles
\[
Q_1=\left[0,\frac12\right]\times\left[0,\frac12\right],\qquad
Q_2=\left[0,\frac12\right]\times\left[\frac12,1\right],
\]
\[
Q_3=\left[\frac12,1\right]\times\left[0,\frac12\right],\qquad
Q_4=\left[\frac12,1\right]\times\left[\frac12,1\right].
\]
On each rectangle \(Q=[a,b]\times[c,d]\), we set
\[
p=a+(b-a)u,\qquad x=c+(d-c)v,\qquad (u,v)\in[0,1]\times[0, 1],
\]
and write the transformed polynomial in the Bernstein basis of bidegree $(6,4)$
\[
F_Q(u,v)=\sum_{i=0}^{6}\sum_{j=0}^{4} b^{(Q)}_{ij} B_i^6(u)B_j^4(v),
\]
where,
\[
B_i^6(u)=\binom{6}{i}u^i(1-u)^{6-i},
\qquad
B_j^4(v)=\binom{4}{j}v^j(1-v)^{4-j}.
\]
Now we list every Bernstein coefficient explicitly.
For $Q_{1},$ where
\[
Q_1=\left[0,\frac12\right]\times\left[0,\frac12\right],
\]
the Bernstein coefficient matrix is
\[
\mathcal B(F,Q_1)=
\begin{pmatrix}
	0 & 0 & \frac{112}{3} & 103 & 196\\[5mm]
	0 & -\frac{34}{3} & \frac{124}{9} & \frac{415}{6} & 158\\[5mm]
	\frac{512}{15} & \frac{29}{3} & \frac{1141}{60} & \frac{3551}{60} & \frac{4123}{30}\\[5mm]
	\frac{199}{2} & \frac{1209}{20} & \frac{2453}{48} & \frac{575}{8} & \frac{5351}{40}\\[5mm]
	\frac{2834}{15} & \frac{64571}{480} & \frac{25151}{240} & \frac{16589}{160} & \frac{17441}{120}\\[5mm]
	\frac{3521}{12} & \frac{7177}{32} & \frac{50297}{288} & \frac{607}{4} & \frac{8261}{48}\\[5mm]
	\frac{25827}{64} & \frac{41391}{128} & \frac{16507}{64} & \frac{27783}{128} & \frac{13983}{64}
\end{pmatrix}.
\]
Thus, we have 
\[
\min \mathcal B(F,Q_1)=-\frac{34}{3}.
\]
For $Q_{2},$ we have
\[
Q_2=\left[0,\frac12\right]\times\left[\frac12,1\right],
\]
the Bernstein coefficient matrix is
\[
\mathcal B(F,Q_2)=
\begin{pmatrix}
	196 & 289 & \frac{1228}{3} & 556 & 736\\[5mm]
	158 & \frac{1481}{6} & \frac{3322}{9} & 528 & 736\\[5mm]
	\frac{4123}{30} & \frac{12941}{60} & \frac{19921}{60} & \frac{4937}{10} & \frac{10774}{15}\\[5mm]
	\frac{5351}{40} & \frac{7827}{40} & \frac{71689}{240} & \frac{3631}{8} & \frac{3414}{5}\\[5mm]
	\frac{17441}{120} & \frac{89761}{480} & \frac{4343}{16} & \frac{198149}{480} & \frac{38131}{60}\\[5mm]
	\frac{8261}{48} & \frac{4619}{24} & \frac{73745}{288} & \frac{12141}{32} & \frac{2353}{4}\\[5mm]
	\frac{13983}{64} & \frac{28149}{128} & \frac{16873}{64} & \frac{47025}{128} & \frac{35631}{64}
\end{pmatrix}.
\]
Thus, we have
\[
\min \mathcal B(F,Q_2)=\frac{5351}{40}.
\]
For $Q_{3}.$ where
\[
Q_3=\left[\frac12,1\right]\times\left[0,\frac12\right],
\]
the Bernstein coefficient matrix is
\[
\mathcal B(F,Q_3)=
\begin{pmatrix}
	\frac{25827}{64} & \frac{41391}{128} & \frac{16507}{64} & \frac{27783}{128} & \frac{13983}{64}\\[5mm]
	\frac{49313}{96} & \frac{27037}{64} & \frac{49133}{144} & \frac{18071}{64} & \frac{25427}{96}\\[5mm]
	\frac{151069}{240} & \frac{7963}{15} & \frac{157649}{360} & \frac{3649}{10} & \frac{26469}{80}\\[5mm]
	\frac{29659}{40} & \frac{2569}{4} & \frac{43627}{80} & \frac{9313}{20} & \frac{4209}{10}\\[5mm]
	\frac{16799}{20} & \frac{90179}{120} & \frac{59713}{90} & \frac{70729}{120} & \frac{10877}{20}\\[5mm]
	\frac{5497}{6} & \frac{10285}{12} & 798 & \frac{8975}{12} & \frac{4295}{6}\\[5mm]
	963 & 963 & 963 & 963 & 963
\end{pmatrix}.
\]
Thus, we have
\[
\min \mathcal B(F,Q_3)=\frac{27783}{128}.
\]
For $Q_{4},$ where
\[
Q_4=\left[\frac12,1\right]\times\left[\frac12,1\right],
\]
the Bernstein coefficient matrix is
\[
\mathcal B(F,Q_4)=
\begin{pmatrix}
	\frac{13983}{64} & \frac{28149}{128} & \frac{16873}{64} & \frac{47025}{128} & \frac{35631}{64}\\[5mm]
	\frac{25427}{96} & \frac{47495}{192} & \frac{2441}{9} & \frac{22743}{64} & \frac{16807}{32}\\[5mm]
	\frac{26469}{80} & \frac{11873}{40} & \frac{21727}{72} & \frac{21883}{60} & \frac{122269}{240}\\[5mm]
	\frac{4209}{10} & \frac{7523}{20} & \frac{29307}{80} & \frac{16367}{40} & \frac{21007}{40}\\[5mm]
	\frac{10877}{20} & \frac{11959}{24} & \frac{21656}{45} & \frac{20291}{40} & \frac{35521}{60}\\[5mm]
	\frac{4295}{6} & \frac{2735}{4} & \frac{2009}{3} & \frac{8191}{12} & \frac{1463}{2}\\[5mm]
	963 & 963 & 963 & 963 & 963
\end{pmatrix}.
\]
Thus, we have 
\[
\min \mathcal B(F,Q_4)=\frac{13983}{64}.
\]
Therefore, after dividing the original square into four equal rectangles, the minimum Bernstein coefficient on each one is
\[
Q_1:\ -\frac{34}{3},\qquad
Q_2:\ \frac{5351}{40},\qquad
Q_3:\ \frac{27783}{128},\qquad
Q_4:\ \frac{13983}{64}.
\]
Therefore Bernstein already proves
\[
F(p_1,x)>0
\]
on \(Q_2\), \(Q_3\), and \(Q_4\), since all Bernstein coefficients there are positive.
The only unresolved rectangle is
\[
Q_1=\left[0,\frac12\right]\times\left[0,\frac12\right],
\]
because there the minimum Bernstein coefficient is negative, precisely
$
-{34}/{3}.
$
So the next step is to subdivide only \(Q_1\) into four smaller rectangles and repeat the Bernstein calculation there.\\[2mm]
Now subdivide
\[
Q_1=\left[0,\frac12\right]\times\left[0,\frac12\right]
\]
into the four smaller rectangles
\[
Q_{11}=\left[0,\frac14\right]\times\left[0,\frac14\right],\qquad
Q_{12}=\left[0,\frac14\right]\times\left[\frac14,\frac12\right],
\]
\[
Q_{13}=\left[\frac14,\frac12\right]\times\left[0,\frac14\right],\qquad
Q_{14}=\left[\frac14,\frac12\right]\times\left[\frac14,\frac12\right].
\]
We continue to work with
\[
F(p_1,x)=1024-G_1(p_1,x),
\qquad p:=p_1.
\]
On each rectangle \(Q=[a,b]\times[c,d]\), set
\[
p=a+(b-a)u,\qquad x=c+(d-c)v,\qquad (u,v)\in[0,1]\times[0, 1],
\]
and write the transformed polynomial in the Bernstein basis of bidegree \((6,4)\):
\[
F_Q(u,v)=\sum_{i=0}^{6}\sum_{j=0}^{4} b^{(Q)}_{ij} B_i^6(u)B_j^4(v),
\]
where
\[
B_i^6(u)=\binom{6}{i}u^i(1-u)^{6-i},
\qquad
B_j^4(v)=\binom{4}{j}v^j(1-v)^{4-j}.
\]
Now we list every Bernstein coefficient explicitly.
For $Q_{11},$ where
\[
Q_{11}=\left[0,\frac14\right]\times\left[0,\frac14\right],
\]
the Bernstein coefficient matrix is\\
\[
\mathcal B(F,Q_{11})=
\begin{pmatrix}
	0 & 0 & \frac{28}{3} & \frac{215}{8} & 52\\[5mm]
	0 & -\frac{17}{6} & \frac{32}{9} & \frac{583}{32} & \frac{163}{4}\\[5mm]
	\frac{128}{15} & \frac{317}{120} & \frac{16567}{2880} & \frac{16421}{960} & \frac{2191}{60}\\[5mm]
	\frac{2019}{80} & \frac{5147}{320} & \frac{119909}{7680} & \frac{119091}{5120} & \frac{12501}{320}\\[5mm]
	\frac{2969}{60} & \frac{566747}{15360} & \frac{750769}{23040} & \frac{1111661}{30720} & \frac{92093}{1920}\\[5mm]
	\frac{30893}{384} & \frac{65853}{1024} & \frac{687673}{12288} & \frac{42439}{768} & \frac{64183}{1024}\\[5mm]
	\frac{480003}{4096} & \frac{1596537}{16384} & \frac{1392057}{16384} & \frac{1307517}{16384} & \frac{338553}{4096}
\end{pmatrix}.
\]
Thus, we have
\[
\min \mathcal B(F,Q_{11})=-\frac{17}{6}.
\]
For $Q_{12}$
\[
Q_{12}=\left[0,\frac14\right]\times\left[\frac14,\frac12\right],
\]
the Bernstein coefficient matrix is\\
\[
\mathcal B(F,Q_{12})=
\begin{pmatrix}
	52 & \frac{617}{8} & \frac{659}{6} & \frac{299}{2} & 196\\[5mm]
	\frac{163}{4} & \frac{2025}{32} & \frac{6745}{72} & \frac{3157}{24} & 177\\[5mm]
	\frac{2191}{60} & \frac{17897}{320} & \frac{240187}{2880} & \frac{18999}{160} & \frac{19483}{120}\\[5mm]
	\frac{12501}{320} & \frac{280941}{5120} & \frac{605459}{7680} & 111 & \frac{48643}{320}\\[5mm]
	\frac{92093}{1920} & \frac{367063}{6144} & \frac{183625}{2304} & \frac{332063}{3072} & \frac{55993}{384}\\[5mm]
	\frac{64183}{1024} & \frac{107671}{1536} & \frac{350787}{4096} & \frac{224847}{2048} & \frac{220651}{1536}\\[5mm]
	\frac{338553}{4096} & \frac{1400907}{16384} & \frac{1578837}{16384} & \frac{1899297}{16384} & \frac{595983}{4096}
\end{pmatrix}.
\]\\
Thus, we have
\[
\min \mathcal B(F,Q_{12})=\frac{2191}{60}.
\]
For $Q_{13},$ where
\[
Q_{13}=\left[\frac14,\frac12\right]\times\left[0,\frac14\right],
\]
the Bernstein coefficient matrix is\\
\[
\mathcal B(F,Q_{13})=
\begin{pmatrix}
	\frac{480003}{4096} & \frac{1596537}{16384} & \frac{1392057}{16384} & \frac{1307517}{16384} & \frac{338553}{4096}\\[5mm]
	\frac{945721}{6144} & \frac{1069713}{8192} & \frac{2800825}{24576} & \frac{2564503}{24576} & \frac{210187}{2048}\\[5mm]
	\frac{1005743}{5120} & \frac{10410323}{61440} & \frac{27388337}{184320} & \frac{8255597}{61440} & \frac{1964059}{15360}\\[5mm]
	\frac{624031}{2560} & \frac{2182609}{10240} & \frac{1925727}{10240} & \frac{346497}{2048} & \frac{404011}{2560}\\[5mm]
	\frac{1132141}{3840} & \frac{4004671}{15360} & \frac{10668157}{46080} & \frac{3196279}{15360} & \frac{736927}{3840}\\[5mm]
	\frac{133817}{384} & \frac{477931}{1536} & \frac{1281499}{4608} & \frac{128205}{512} & \frac{29395}{128}\\[5mm]
	\frac{25827}{64} & \frac{93045}{256} & \frac{83725}{256} & \frac{75663}{256} & \frac{17325}{64}
\end{pmatrix}.
\]\\
Thus, we have 
\[
\min \mathcal B(F,Q_{13})=\frac{1307517}{16384}.
\]
For $Q_{14},$ where
\begin{align*}
	Q_{14}=\left[\frac14,\frac12\right]\times\left[\frac14,\frac12\right],
\end{align*}\\
the Bernstein coefficient matrix is\\
\[
\mathcal B(F,Q_{14})=
\begin{pmatrix}
	\frac{338553}{4096} & \frac{1400907}{16384} & \frac{1578837}{16384} & \frac{1899297}{16384} & \frac{595983}{4096}\\[5mm]
	\frac{210187}{2048} & \frac{2479985}{24576} & \frac{877263}{8192} & \frac{999909}{8192} & \frac{905345}{6144}\\[5mm]
	\frac{1964059}{15360} & \frac{497125}{4096} & \frac{4519201}{36864} & \frac{1629815}{12288} & \frac{156895}{1024}\\[5mm]
	\frac{404011}{2560} & \frac{1499603}{10240} & \frac{1459963}{10240} & \frac{302057}{2048} & \frac{417471}{2560}\\[5mm]
	\frac{736927}{3840} & \frac{2699137}{15360} & \frac{1537061}{9216} & \frac{852653}{5120} & \frac{226571}{1280}\\[5mm]
	\frac{29395}{128} & \frac{106955}{512} & \frac{898999}{4608} & \frac{291607}{1536} & \frac{74993}{384}\\[5mm]
	\frac{17325}{64} & \frac{62937}{256} & \frac{58273}{256} & \frac{55749}{256} & \frac{13983}{64}
\end{pmatrix}.
\]\\
Thus, we have 
\[
\min \mathcal B(F,Q_{14})=\frac{338553}{4096}.
\]
Therefore, after subdividing \(Q_1\) into four smaller rectangles, the minimum Bernstein coefficient on each one is
\[
Q_{11}:\ -\frac{17}{6},\qquad
Q_{12}:\ \frac{2191}{60},\qquad
Q_{13}:\ \frac{1307517}{16384},\qquad
Q_{14}:\ \frac{338553}{4096}.
\]
So Bernstein now proves
\[
F(p_1,x)>0
\]
on \(Q_{12}\), \(Q_{13}\), and \(Q_{14}\), since all Bernstein coefficients there are positive.
The only unresolved rectangle after this second subdivision is
\[
Q_{11}=\left[0,\frac14\right]\times\left[0,\frac14\right],
\]
because there the minimum Bernstein coefficient is still negative, precisely
$
-{17}/{6}.
$
So the next step is to subdivide only \(Q_{11}\) again into four smaller rectangles and repeat the Bernstein calculation there.\\[2mm]
Now subdivide
\[
Q_{11}=\left[0,\frac14\right]\times\left[0,\frac14\right]
\]
into the four smaller rectangles
\[
Q_{111}=\left[0,\frac18\right]\times\left[0,\frac18\right],\qquad
Q_{112}=\left[0,\frac18\right]\times\left[\frac18,\frac14\right],
\]
\[
Q_{113}=\left[\frac18,\frac14\right]\times\left[0,\frac18\right],\qquad
Q_{114}=\left[\frac18,\frac14\right]\times\left[\frac18,\frac14\right].
\]
We still work with
\[
F(p,x)=1024-G_1(p,x),\qquad p:=p_1.
\]
For each rectangle \(Q=[a,b]\times[c,d]\), set
\[
p=a+(b-a)u,\qquad x=c+(d-c)v,\qquad (u,v)\in[0,1]\times[0, 1],
\]
and write the transformed polynomial in the Bernstein basis of bidegree \((6,4)\),
\[
F_Q(u,v)=\sum_{i=0}^{6}\sum_{j=0}^{4} b_{ij}^{(Q)}\,B_i^6(u)B_j^4(v),
\]
where
\[
B_i^6(u)=\binom{6}{i}u^i(1-u)^{6-i},
\qquad
B_j^4(v)=\binom{4}{j}v^j(1-v)^{4-j}.
\]
Now the Bernstein coefficient matrices are as follows.
For $Q_{111},$ where
\begin{align*}
	Q_{111}=\left[0,\frac18\right]\times\left[0,\frac18\right],
\end{align*}\\
\[
\mathcal B(F,Q_{111})=
\begin{pmatrix}
	0 & 0 & \frac{7}{3} & \frac{439}{64} & \frac{431}{32}\\[5mm]
	0 & -\frac{17}{24} & \frac{65}{72} & \frac{7225}{1536} & \frac{2713}{256}\\[5mm]
	\frac{32}{15} & \frac{661}{960} & \frac{23693}{15360} & \frac{2197}{480} & \frac{59645}{6144}\\[5mm]
	\frac{4067}{640} & \frac{21231}{5120} & \frac{1034701}{245760} & \frac{2109869}{327680} & \frac{88153}{8192}\\[5mm]
	\frac{6049}{480} & \frac{944287}{98304} & \frac{2174381}{245760} & \frac{13412569}{1310720} & \frac{26941817}{1966080}\\[5mm]
	\frac{255581}{12288} & \frac{556783}{32768} & \frac{72585779}{4718592} & \frac{16662671}{1048576} & \frac{58100225}{3145728}\\[5mm]
	\frac{8089155}{262144} & \frac{54984213}{2097152} & \frac{99545727}{4194304} & \frac{195731055}{8388608} & \frac{209583495}{8388608}
\end{pmatrix}.
\]
Thus, we have
\[
\min \mathcal B(F,Q_{111})=-\frac{17}{24}.
\]
For $Q_{112},$ where
\begin{align*}
	Q_{112}=\left[0,\frac18\right]\times\left[\frac18,\frac14\right],
\end{align*}\\
\[
\mathcal B(F,Q_{112})=
\begin{pmatrix}
	\frac{431}{32} & \frac{1285}{64} & \frac{1381}{48} & \frac{631}{16} & 52\\[5mm]
	\frac{2713}{256} & \frac{25331}{1536} & \frac{28199}{1152} & \frac{4411}{128} & \frac{371}{8}\\[5mm]
	\frac{59645}{6144} & \frac{227921}{15360} & \frac{338927}{15360} & \frac{80139}{2560} & \frac{10201}{240}\\[5mm]
	\frac{88153}{8192} & \frac{4942371}{327680} & \frac{2641727}{122880} & \frac{2452339}{81920} & \frac{103317}{2560}\\[5mm]
	\frac{26941817}{1966080} & \frac{67529561}{3932160} & \frac{7447817}{327680} & \frac{29796421}{983040} & \frac{1225589}{30720}\\[5mm]
	\frac{58100225}{3145728} & \frac{66212437}{3145728} & \frac{121259051}{4718592} & \frac{4243233}{131072} & \frac{1345615}{32768}\\[5mm]
	\frac{209583495}{8388608} & \frac{223435935}{8388608} & \frac{127250607}{4194304} & \frac{75626793}{2097152} & \frac{11484225}{262144}
\end{pmatrix}.
\]
Thus, we have
\[
\min \mathcal B(F,Q_{112})=\frac{59645}{6144}.
\]
For $Q_{113},$ where
\begin{align*}
	Q_{113}=\left[\frac18,\frac14\right]\times\left[0,\frac18\right],
\end{align*}
\[
\mathcal B(F,Q_{113})=
\begin{pmatrix}
	\frac{8089155}{262144} & \frac{54984213}{2097152} & \frac{99545727}{4194304} & \frac{195731055}{8388608} & \frac{209583495}{8388608}\\[5mm]
	\frac{16088873}{393216} & \frac{37167157}{1048576} & \frac{605568427}{18874368} & \frac{129080371}{4194304} & \frac{396349585}{12582912}\\[5mm]
	\frac{51939757}{983040} & \frac{73158895}{1572864} & \frac{1993607707}{47185920} & \frac{419448987}{10485760} & \frac{416937499}{10485760}\\[5mm]
	\frac{2180003}{32768} & \frac{77773049}{1310720} & \frac{28389403}{524288} & \frac{267038767}{5242880} & \frac{260564127}{5242880}\\[5mm]
	\frac{6708687}{81920} & \frac{145129691}{1966080} & \frac{798453347}{11796480} & \frac{499235521}{7864320} & \frac{160493139}{2621440}\\[5mm]
	\frac{2428585}{24576} & \frac{17664895}{196608} & \frac{32542313}{393216} & \frac{60985241}{786432} & \frac{58387873}{786432}\\[5mm]
	\frac{480003}{4096} & \frac{3516549}{32768} & \frac{6505143}{65536} & \frac{12193311}{131072} & \frac{11621391}{131072}
\end{pmatrix}.
\]
Thus, we have
\[
\min \mathcal B(F,Q_{113})=\frac{195731055}{8388608}.
\]
For $Q_{114},$ where
\[
Q_{114}=\left[\frac18,\frac14\right]\times\left[\frac18,\frac14\right],
\]
\[
\mathcal B(F,Q_{114})=
\begin{pmatrix}
	\frac{209583495}{8388608} & \frac{223435935}{8388608} & \frac{127250607}{4194304} & \frac{75626793}{2097152} & \frac{11484225}{262144}\\[5mm]
	\frac{396349585}{12582912} & \frac{405458057}{12582912} & \frac{660219259}{18874368} & \frac{41680929}{1048576} & \frac{6101765}{131072}\\[5mm]
	\frac{416937499}{10485760} & \frac{414426011}{10485760} & \frac{1948400923}{47185920} & \frac{354397343}{7864320} & \frac{50008423}{983040}\\[5mm]
	\frac{260564127}{5242880} & \frac{254089487}{5242880} & \frac{25799547}{524288} & \frac{68079949}{1310720} & \frac{1858157}{32768}\\[5mm]
	\frac{160493139}{2621440} & \frac{463723313}{7864320} & \frac{691916723}{11796480} & \frac{39525309}{655360} & \frac{15727231}{245760}\\[5mm]
	\frac{58387873}{786432} & \frac{18596835}{262144} & \frac{9115859}{131072} & \frac{13782067}{196608} & \frac{595285}{8192}\\[5mm]
	\frac{11621391}{131072} & \frac{11049471}{131072} & \frac{5361303}{65536} & \frac{2661729}{32768} & \frac{338553}{4096}
\end{pmatrix}.
\]
Thus, we have
\[
\min \mathcal B(F,Q_{114})=\frac{209583495}{8388608}.
\]
Therefore, after subdividing \(Q_{11}\) into four smaller rectangles, the minimum Bernstein coefficient on each one is
\[
Q_{111}:\ -\frac{17}{24},\qquad
Q_{112}:\ \frac{59645}{6144},\qquad
Q_{113}:\ \frac{195731055}{8388608},\qquad
Q_{114}:\ \frac{209583495}{8388608}.
\]
So Bernstein proves
\[
F(p_1,x)>0
\]
on \(Q_{112}\), \(Q_{113}\), and \(Q_{114}\).
The only unresolved rectangle after this third subdivision is
\[
Q_{111}=\left[0,\frac18\right]\times\left[0,\frac18\right],
\]
because there the minimum Bernstein coefficient is still negative, precisely
$
-{17}/{24}.
$
So at this stage, subdivision plus Bernstein settles every part except the tiny corner square \(Q_{111}\). That last square must be handled by a direct local estimate for \(F=1024-G_1\).\\[2mm]
On the last square
\[
Q_{111}=\left[0,\frac18\right]\times\left[0,\frac18\right],
\]
set
\[
F(p,x):=1024-G_1(p,x), \qquad p:=p_1.
\]
We prove directly that
\[
F(p,x)\ge 0 \qquad \text{for } (p,x)\in Q_{111},
\]
and in fact
\[
F(p,x)>0 \qquad \text{for } (p,x)\neq(0,0).
\]
First expand $F$ such that 
\[
\begin{aligned}
	F(p,x)
	&=-61p^6+464p^5-1024p^4-464p^3+2048p^2 \\
	&\quad+\left(244p^6+640p^5+620p^4+448p^3-864p^2-1088p\right)x \\
	&\quad+\left(-248p^6-720p^5+1856p^4+976p^3-2504p^2-256p+896\right)x^2 \\
	&\quad+\left(64p^6-640p^5-416p^4-448p^3+640p^2+1088p-288\right)x^3 \\
	&\quad+\left(-128p^6+256p^5+384p^4-512p^3-384p^2+256p+128\right)x^4.
\end{aligned}
\]
Now separate the quadratic part of $p$
\[
F(p,x)=Q(p,x)+R(p,x),
\]
where
\[
Q(p,x)=2048p^2-1088px+896x^2,
\]
and
\[
\begin{aligned}
	R(p,x)
	&=-61p^6+464p^5-1024p^4-464p^3 \\
	&\quad+\left(244p^6+640p^5+620p^4+448p^3-864p^2\right)x \\
	&\quad+\left(-248p^6-720p^5+1856p^4+976p^3-2504p^2-256p\right)x^2 \\
	&\quad+\left(64p^6-640p^5-416p^4-448p^3+640p^2+1088p-288\right)x^3 \\
	&\quad+\left(-128p^6+256p^5+384p^4-512p^3-384p^2+256p+128\right)x^4.
\end{aligned}
\]
Every monomial in \(R\) has total degree at least \(3\).
Now estimate \(Q\) from below. Since
\[
2px\le p^2+x^2,
\]
we get
\[
1088px\le 544p^2+544x^2.
\]
Therefore
\[
Q(p,x)\ge (2048-544)p^2+(896-544)x^2
=1504p^2+352x^2
\ge 352(p^2+x^2).
\]
Next estimate \(R\) on \(Q_{111}\), where
\[
0\le p\le \frac18,\qquad 0\le x\le \frac18.
\]
Write
\[
R(p,x)=\sum_{i+j\ge 3} c_{ij}p^ix^j.
\]
For each monomial \(p^ix^j\) with \(i+j\ge 3\), we have
\[
p^ix^j\le 8^{-(i+j-2)}(p^2+x^2).
\]
Indeed:
if $i \ge 2$, then
\[
p^ix^j=p^2p^{\,i-2}x^j\le 8^{-(i+j-2)}p^2\le 8^{-(i+j-2)}(p^2+x^2),
\]
if $i=1$, then $j\ge 2$, so
\[
p x^j=x^2(px^{j-2})\le 8^{-(j-1)}x^2=8^{-(i+j-2)}x^2\le 8^{-(i+j-2)}(p^2+x^2),
\]
if $i=0$, then $j\ge 3$, so
\[
x^j\le 8^{-(j-2)}x^2=8^{-(i+j-2)}x^2\le 8^{-(i+j-2)}(p^2+x^2).
\]
Hence
\[
|R(p,x)|
\le
\left(\sum_{i+j\ge 3}|c_{ij}|\,8^{-(i+j-2)}\right)(p^2+x^2).
\]
Now compute that coefficient sum explicitly from the coefficients of \(R\):
\[
\begin{aligned}
	\sum_{i+j\ge 3}|c_{ij}|\,8^{-(i+j-2)}
	&=
	\frac{61}{4096}
	+\frac{464}{512}
	+\frac{1024}{64}
	+\frac{464}{8}
	+\frac{244}{512}
	+\frac{640}{64}
	+\frac{620}{8}
	+\frac{448}{1}
	+\frac{864}{8}
	\\
	&\quad
	+\frac{248}{64}
	+\frac{720}{8}
	+\frac{1856}{1}
	+\frac{976}{1}
	+\frac{2504}{8}
	+\frac{256}{1}
	+\frac{64}{8}
	+\frac{640}{1}
	+\frac{416}{1}
	+\frac{448}{1}
	\\
	&\quad
	+\frac{640}{8}
	+\frac{1088}{1}
	+288
	+128
	+\frac{256}{8}
	+\frac{384}{1}
	+\frac{512}{1}
	+\frac{384}{8}
	+256
	+128
	\\
	&=\frac{42177473}{131072}.
\end{aligned}
\]
Therefore
\[
|R(p,x)|\le \frac{42177473}{131072}(p^2+x^2).
\]
Combining the two estimates,
\[
\begin{aligned}
	F(p,x)
	&=Q(p,x)+R(p,x)\\
	&\ge 352(p^2+x^2)-\frac{42177473}{131072}(p^2+x^2)\\
	&=\frac{3959871}{131072}(p^2+x^2).
\end{aligned}
\]
Since
\[
\frac{3959871}{131072}>0,
\]
we conclude that
\[
F(p,x)\ge 0 \qquad \text{on }Q_{111},
\]
and in fact
\[
F(p,x)>0 \qquad \text{for }(p,x)\neq(0,0).
\]
Thus
\[
1024-G_1(p_1,x)\ge 0
\qquad\text{on}\qquad
\left[0,\frac18\right]\times\left[0,\frac18\right].
\]
Together with the Bernstein positivity on all the other subrectangles, this completes the proof that
\[
1024-G_1(p_1,x)\ge 0
\qquad\text{for all }(p_1,x)\in[0,1]\times[0, 1],
\]
with equality only at
\[
(p_1,x)=(0,0).
\]
So finally,
\[
G_1(p_1,x)\le 1024
\quad\text{on}\quad [0,1]\times[0, 1],
\]
and the maximum is \(1024\), attained only at \((0,0)\).\\[2mm]
Now we deal with the second function $G_{2}(p_1, x)$.\\[2mm]
Let
\[
F(p_1,x):=G_2(p_1,x),
\qquad 0\le p_1\le 1,\quad 0\le x\le 1.
\]
For simplicity, write $p=p_{1}$. Then
\[
\begin{aligned}
	F(p,x)=\,&61 p^6 + 244 p^4 (1-p^2)x + 8p^2(89-120p^2+31p^4)x^2 \\
	&-32(-9-7p^2+14p^4+2p^6)x^3 + 128p^2(1-p^2)^2x^4 \\
	&+288(1-p^2)\bigl(3p^2+4(1-p^2)x\bigr)(1-x^2) \\
	&+16p(1-p^2)(1-x^2)\bigl(29p^2+(68+40p^2)x+16(1-p^2)x^2\bigr).
\end{aligned}
\]
After expanding and collecting like terms, we obtain
\[
\begin{aligned}
	F(p,x)=\,&128p^{6}x^{4}-64p^{6}x^{3}+248p^{6}x^{2}-244p^{6}x+61p^{6} \\
	&-256p^{5}x^{4}+640p^{5}x^{3}+720p^{5}x^{2}-640p^{5}x-464p^{5} \\
	&-256p^{4}x^{4}-1600p^{4}x^{3}-96p^{4}x^{2}+1396p^{4}x-864p^{4} \\
	&+512p^{3}x^{4}+448p^{3}x^{3}-976p^{3}x^{2}-448p^{3}x+464p^{3} \\
	&+128p^{2}x^{4}+2528p^{2}x^{3}-152p^{2}x^{2}-2304p^{2}x+864p^{2} \\
	&-256px^{4}-1088px^{3}+256px^{2}+1088px \\
	&-864x^{3}+1152x.
\end{aligned}
\]
This is a polynomial of bidegree $(6,4)$. Hence on the unit square $[0,1]\times[0, 1]$ it can be written in the Bernstein basis as
\[
F(p,x)=\sum_{i=0}^{6}\sum_{j=0}^{4} b_{ij}\,B_i^6(p)\,B_j^4(x),
\]
where
\[
B_i^6(p)=\binom{6}{i}p^i(1-p)^{6-i},
\qquad
B_j^4(x)=\binom{4}{j}x^j(1-x)^{4-j}.
\]
For the whole rectangle \([0,1]\times[0,1]\), the Bernstein coefficients are\\
\[
\bigl(b_{ij}\bigr)=
\begin{pmatrix}
	0 & 288 & 576 & 648 & 288\\[5mm]
	0 & \dfrac{1000}{3} & \dfrac{6064}{9} & 760 & 288\\[5mm]
	\dfrac{288}{5} & \dfrac{5968}{15} & \dfrac{2252}{3} & \dfrac{12772}{15} & \dfrac{5384}{15}\\[5mm]
	196 & \dfrac{2496}{5} & \dfrac{12158}{15} & 910 & \dfrac{2504}{5}\\[5mm]
	\dfrac{1904}{5} & \dfrac{3103}{5} & \dfrac{7606}{9} & \dfrac{13583}{15} & \dfrac{9284}{15}\\[5mm]
	\dfrac{1328}{3} & 607 & \dfrac{2170}{3} & \dfrac{2159}{3} & 524\\[5mm]
	61 & 61 & 61 & 61 & 61
\end{pmatrix}.
\]\\
Now we use the basic Bernstein-basis enclosure property: if
\[
F(p,x)=\sum_{i=0}^{6}\sum_{j=0}^{4} b_{ij}\,B_i^6(p)\,B_j^4(x)
\qquad \text{on }\quad  [0,1]\times[0, 1],
\]
then
\[
\min_{i,j} b_{ij}\le F(p,x)\le \max_{i,j} b_{ij}
\qquad \text{for all } (p,x)\in[0,1]\times[0, 1].
\]
From the coefficient matrix above, the largest Bernstein coefficient is
\[
\max_{0\le i\le 6,\,0\le j\le 4} b_{ij}=910.
\]
Therefore,
\[
F(p,x)\le 910
\qquad \text{for all } (p,x)\in[0,1]\times[0, 1].
\]
Equivalently,
\begin{align*}
	G_2(p_1,x)\le 910
	\qquad \text{for all } 0\le p_1\le 1,\,0\le x\le 1.
\end{align*}
Thus we obtain the desired upper bound
\[
\max_{0\le p_1\le 1,\,0\le x\le 1} G_2(p_1,x)\le 910.
\]
Combining both cases we obtain that 
\[
H_{1}(p_1, x, y)\le 1024 \quad \text{over} \quad \mathcal{D}.
\]
Therefore 
\[
|H_{3}(1)|\le \frac{1024}{9216}=\frac{1}{9}.
\]
Finally, the estimate is sharp. Taking $w(z)=z^3$ in the defining subordination,
\[
\frac{z f'(z)}{f(z)}=\left(1+\frac{z^3}{2}\right)^2,
\]
and integrating gives
\[
f(z)=z\exp\!\left(\int_0^z \frac{(1+t^3/2)^2-1}{t}\,dt\right)
=z\exp\!\left(\frac{z^3}{3}+\frac{z^6}{24}\right).
\]
For this function, $a_2=a_3=a_5=0$ and $a_4=1/3$, hence
\[
H_3(1)=-a_4^2=-\frac19,
\]
Therefore, we have $|H_3(1)|=1/9$. This completes the proof.
\end{proof}

\end{document}
